%% =====================================================================
\section{Comparative Analysis}
\label{sec:comparative}
%% =====================================================================

\subsection{Capability Coverage}
CacheCollapse integrates concerns that prior work addresses only in prt
(Table~\ref{tab:capability}). Performance-oriented stampede work
\cite{vattani2015,nishtala2013} contributes coalescing and early recomputation but no
attacker model, detection, or testbed. Application-layer DDoS work
\cite{praseed2019,xie2009,ranjan2009} contributes detection and scheduling but neither
the cache-specific mechanism nor a reproducible container environment. Work on
cache- and CDN-targeting attacks \cite{nguyen2019,guo2020,triukose2009} demonstrates
attacks on shared edge caches but does not address the application-internal cache or
layered mitigation. Container-security work \cite{combe2016,sultan2019} catalogs
weaknesses without instrumenting a specific availability attack end to end. The Yo-Yo
attack \cite{sides2016,bremlerbarr2017} is closest in spirit but targets elasticity
rather than caching. Overload-control work \cite{welsh2001,welsh2003} supplies general
mechanisms without an attack model.

\begin{table}[t]
\centering
\caption{Capability coverage. \yes~= addressed; \prt~= partial; \no~= not addressed.
``Analytical'' denotes a quantitative model of the attack or its mitigation.}
\label{tab:capability}
\small
\begin{tabular}{@{}lcccccc@{}}
\toprule
Work & Adversarial & Testbed & Analytical & Detection & Layered mit. & Containerized \\
\midrule
Probabilistic prevention \cite{vattani2015} & \no & \no & \yes & \no & \prt & \no \\
Memcache leases \cite{nishtala2013} & \no & \no & \no & \no & \prt & \no \\
App-layer DDoS \cite{xie2009,praseed2019} & \yes & \no & \prt & \yes & \no & \no \\
DDoS-Shield \cite{ranjan2009} & \yes & \prt & \yes & \yes & \prt & \no \\
Flash crowds \cite{jung2002,kandula2005} & \yes & \prt & \no & \yes & \prt & \no \\
DDoS surveys \cite{mirkovic2004,zargar2013} & \yes & \no & \no & \prt & \prt & \no \\
CPDoS / CDN abuse \cite{nguyen2019,guo2020} & \yes & \yes & \no & \no & \prt & \no \\
Container security \cite{combe2016,sultan2019} & \prt & \no & \no & \no & \no & \prt \\
Yo-Yo attack \cite{sides2016,bremlerbarr2017} & \yes & \prt & \yes & \prt & \no & \yes \\
Overload control \cite{welsh2001,welsh2003} & \no & \yes & \prt & \prt & \prt & \no \\
\textbf{CacheCollapse} & \yes & \yes & \yes & \yes & \yes & \yes \\
\bottomrule
\end{tabular}
\end{table}

\subsection{Comparison of Mitigation Techniques}
Table~\ref{tab:mitigations} compares the mitigations available against a cache stampede
along six dimensions. Two observations stand out. First, the techniques fall into two
groups: those that reduce the \emph{probability} of a synchronized miss (TTL jitter,
early recomputation, longer TTLs) and those that bound its \emph{cost} once it happens
(single-flight, stale serving, rate limiting, admission control). An adversary who
can time bursts defeats the first group---jitter adds noise the attacker can average
away by probing, and early recomputation only refreshes keys that are being read
anyway---so a security-oriented design must include the second group. Second, within
the second group only edge-level controls protect the proxy and connection path. This
is the empirical result of Section~\ref{sec:ablation} restated as a design rule.

\begin{table}[t]
\centering
\caption{Comparison of stampede mitigations. ``Root cause'' = removes redundant
regeneration; ``Proxy'' = protects the proxy/connection path; ``Adversarial'' = remains
effective when bursts are deliberately timed.}
\label{tab:mitigations}
\footnotesize
\begin{tabularx}{\linewidth}{@{}L{3.1cm}cccccY@{}}
\toprule
Technique & Root cause & Protects DB & Proxy & Adversarial & Staleness & Cost / caveat \\
\midrule
Longer TTL & \no & \prt & \no & \no & Higher & Fewer but equally large stampedes (Eq.~\ref{eq:stampede}) \\
TTL jitter & \no & \prt & \no & \no & Same & De-synchronizes bulk fills only \\
Early recomputation (XFetch) \cite{vattani2015} & \prt & \prt & \no & \prt & Same & Needs regeneration-cost estimate; read-triggered \\
Single-flight lock / lease \cite{nishtala2013,singleflight} & \yes & \yes & \no & \yes & Same & Losers must wait or be served something \\
Stale-while-revalidate \cite{rfc5861} & \prt & \yes & \no & \yes & Bounded & Requires a stale copy (pre-warm) \\
App-layer rate limit & \no & \prt & \no & \prt & None & Acts after the proxy and app have paid for the request \\
Edge rate limit \cite{nginxlimitreq} & \no & \prt & \yes & \yes & None & Per-address; weak against many sources \\
Admission control \cite{welsh2003,cherkasova2002} & \no & \yes & \prt & \yes & None & Needs a load signal; may reject legitimate users \\
Auto-scaling & \no & \prt & \prt & \no & None & Slow; itself attackable (Yo-Yo) \cite{sides2016} \\
\textbf{CacheCollapse (lock + SWR + app + edge RL)} & \yes & \yes & \yes & \yes & Bounded (60 s) & Per-address limit; see Sec.~\ref{sec:limitations} \\
\bottomrule
\end{tabularx}
\end{table}

\subsection{Comparison of Detection Approaches}
Table~\ref{tab:detcompare} positions the two CacheCollapse detectors relative to
generic approaches. The comparison underlines why we chose to deploy both: a
volume threshold is cheap but cannot tell harm from popularity
(Table~\ref{tab:single}); behavioral models of browsing \cite{xie2009} and
session-suspicion scoring \cite{ranjan2009} require per-client modeling that a
stampede, made of identical requests for public objects, does not reward; supervised
classifiers require labelled attacks that rarely exist in production. The rule and
the unsupervised forest need neither per-client state nor labels.

\begin{table}[t]
\centering
\caption{Detection approaches for the cache-stampede IDDoS.}
\label{tab:detcompare}
\footnotesize
\begin{tabularx}{\linewidth}{@{}L{3.2cm}L{1.7cm}L{1.8cm}L{1.8cm}Y@{}}
\toprule
Approach & Labels needed & Detects harm or novelty & Explainable & Main weakness here \\
\midrule
Request-volume threshold & No & Novelty (volume) & Yes & Fires on absorbed bursts (97\%) \\
Per-client browsing model \cite{xie2009} & Normal only & Novelty & Partly & Stampede requests are individually normal \\
Session suspicion scoring \cite{ranjan2009} & No & Novelty & Partly & Many sources defeat per-session state \\
Supervised classifier & Attack + normal & Either & Rarely & Labelled attacks unavailable \\
\textbf{Multi-signal rule (ours)} & No & Harm & Fully & Conservative thresholds; recall 0.27 \\
\textbf{IsolationForest (ours)} & Normal only & Novelty & Weakly & Needs clean window; flags idleness \\
\textbf{Rule OR forest (ours)} & Normal only & Both & Partly & 12/12 episodes, FPR 0.3\% \\
\bottomrule
\end{tabularx}
\end{table}

\subsection{Positioning}
In summary, CacheCollapse is not a new mitigation primitive: single-flight
regeneration, stale serving, and rate limiting all predate it. Its contribution is to
treat the stampede as an adversarial availability problem and to show,
quantitatively and reproducibly, how these primitives must be composed, what each one
protects, how quickly they act, and how the attack can be detected without labels and
without confusing harm with popularity.
